\section*{Hoofdstuk \ref{ch:b2:guided-waves} --- Geleide golven en trilholten}

\begin{solution}{exo:b2:guided-waves:1}
$f_n = nc/2a = 7.5$, $15$, \qty{22.5}{GHz}. Bij \qty{12}{GHz} alleen
$n = 1$: $k_g =
(2\pi/c)\sqrt{f^2 - f_1^2} = \qty{196}{rad/m}$,
$\lambda_g = \qty{3.2}{cm}$, $v_\varphi = c/0.78
= \qty{3.8e8}{m/s}$,
$v_g = \qty{2.3e8}{m/s}$. Bij \qty{6}{GHz}:
$\ell = c/2\pi\sqrt{f_1^2 -
f^2} = \qty{1.1}{cm}$.
\end{solution}

\begin{solution}{exo:b2:guided-waves:2}
TE$_{10}$ \qty{6.55}{GHz}, TE$_{20}$ \qty{13.1}{GHz}, TE$_{01}$
\qty{14.7}{GHz}: enkelmodig van $6.55$ tot \qty{13.1}{GHz}. Bij
\qty{10}{GHz}: $\sqrt{1 - (6.55/10)^2} =
0.756$,
$\lambda_g = \qty{4.0}{cm}$, $v_g = 0.76c$, $\cos\theta = 0.756$,
$\theta = \qty{41}{\degree}$. Met $b \approx a/2$ blijft TE$_{01}$ boven
TE$_{20}$ (de breedst mogelijke enkelmodige band) terwijl er zo veel
mogelijk ruimte voor het veld overblijft (vermogen vóór doorslag).
\end{solution}

\begin{solution}{exo:b2:guided-waves:3}
$f = 1.5 \times 10^8\sqrt{(m/0.3)^2 + (n/0.3)^2 + (p/0.2)^2}$:
\qty{0.71}{GHz}, \qty{0.90}{GHz}, \qty{1.35}{GHz}, \qty{1.95}{GHz}.
$N = 8\pi Vf^3/3c^3 = 8\pi \times 0.018 \times 1.56 \times 10^{28}/
8.1 \times 10^{25} \approx 260$.
Kubus: $f = c\sqrt2/2L$, $L = \qty{2.3}{cm}$.
\end{solution}

\begin{solution}{exo:b2:guided-waves:4}
$\mathrm{NA} = \sqrt{1.48^2 - 1.46^2} = 0.24$, \qty{14}{\degree};
$\Delta t = (n_1L/c)(n_1/n_2
- 1) = \qty{68}{ns/km}$; \qty{135}{ns} over
\qty{2}{km}: \qty{7}{Mbit/s}; met $0.003$: \qty{10}{ns/km},
\qty{50}{Mbit/s}.
\end{solution}

\begin{solution}{exo:b2:guided-waves:5}
(a) $\underline B_y = (k_g/\omega)\underline E$,
$\underline B_z = (\iu/\omega)\partial_y\underline E = (\iu\pi/a\omega)
E_0\cos(\pi y/a)\eu^{\iu(\omega t - k_gz)}$:
in kwadratuur. (b) In $y = 0$ is $B_y = 0$ en
$\vect j_s = \vect e_y\wedge\vect B/\mu_0 = (B_z/\mu_0)\vect e_x$: langs het
veld. (c) $\langle\Pi_z\rangle = E_0^2k_g\sin^2(\pi y/a)/2\mu_0\omega$,
vermogen per eenheid breedte $E_0^2k_ga/4\mu_0\omega$;
$\langle\Pi_y\rangle \propto \operatorname{Re}(\underline E\,\iu\underline B_z^*)
= 0$.
(d) Gemiddelde energie per lengte- en breedte-eenheid
$\tfrac14\varepsilon_0E_0^2a$ (elektrische en magnetische helft elk
$\tfrac18\varepsilon_0E_0^2a$); de verhouding is $c^2k_g/\omega = v_g$.
\end{solution}

\begin{solution}{exo:b2:guided-waves:6}
(a) $f_1 = \qty{150}{GHz} \gg f$: $\ell \approx a/\pi = \qty{0.32}{mm}$;
$\eu^{-0.5/0.32}$: \qty{-14}{dB} in amplitude, $-27$ in vermogen. (b)
$f_1 = \qty{750}{MHz}$; bij \qty{100}{MHz} is $\ell = \qty{6.4}{cm}$:
\qty{-136}{dB} over een meter. (c) $f_1 = \qty{15}{GHz}$,
$\ell =
\qty{3.2}{mm}$: \qty{27}{dB/cm}. (d) De golf past er niet in;
evanescentie is meetkunde, geen dissipatie.
\end{solution}

\begin{solution}{exo:b2:guided-waves:7}
(a)
$P = \tfrac14\varepsilon_0cE_0^2ab\sqrt{1 - f_{10}^2/f^2} = \qty{120}{kW}$.
(b) $\times9$: \qty{1}{MW}; stikstof onder druk of SF$_6$ verhogen het
doorslagveld. (c) $j_s \approx E_0/\mu_0c = \qty{2.7}{kA/m}$;
$\tfrac12R_sj_s^2 = \qty{9e4}{W/m^2}$ over $2a =
\qty{0.046}{m^2}$ per
meter: \qty{4}{kW/m}, $3.5\%$ per meter, ongeveer \qty{0.15}{dB/m}. (d) Drie
keer minder dan de kabel, en megawatten in plaats van kilowatten.
\end{solution}

\begin{solution}{exo:b2:guided-waves:8}
(a) $f = c\sqrt2/2L = \qty{2.12}{GHz}$. (b) $\delta = \qty{1.4}{\micro m}$,
$R_s = \qty{12}{m\ohm}$; $W = \varepsilon_0E_0^2V/8 = \qty{1.1e-15}{}E_0^2$;
$P = \tfrac12R_s(E_0/377)^2 \times 6L^2 =
\qty{2.5e-9}{}E_0^2$:
$Q = \omega W/P \approx 6000$. (c) \qty{360}{kHz}; \qty{0.45}{\micro s}. (d)
$Q \sim 10^{10}$: een veld van tientallen megavolt per meter wordt dan met
kilowatten in plaats van gigawatten onderhouden.
\end{solution}

\begin{solution}{exo:b2:guided-waves:9}
(a) De punten $(m, n, p)$ met $m^2 + n^2 + p^2 \le (2fL/c)^2$ vullen een
achtste bol: $\tfrac18\cdot\tfrac43\pi(2fL/c)^3$, maal twee polarisaties:
$8\pi L^3f^3/3c^3$. (b) $8\pi f^2/c^3$. (c) $9 \times 10^{-7}$ per hertz
(één per megahertz) bij \qty{1}{GHz}; $2 \times 10^5$ per hertz bij
\qty{500}{THz}. (d) Modes zijn discreet zolang hun onderlinge afstand hun
breedte overtreft; in de oven is dat ongeveer één mode per \qty{10}{MHz}
rond \qty{2.45}{GHz} --- tientallen binnen enkele procenten.
\end{solution}

\begin{solution}{exo:b2:guided-waves:10}
(a) Uit $\operatorname{\vect{curl}}\vect E = -\partial_t\vect B$:
\[
\underline B_x = -\frac{k_g}\omega\underline E , \qquad
\underline B_z = \frac{\iu\pi}{a\omega}E_0\cos\frac{\pi x}a\,\eu^{\iu(\omega t - k_gz)} ,
\]
en
$\partial_xB_x + \partial_zB_z = -(k_g\pi/a\omega)E_0\cos(\pi x/a) + (k_g\pi/a\omega)E_0
\cos(\pi x/a) = 0$.
(b)
$(\operatorname{\vect{curl}}\vect B)_y = \partial_zB_x - \partial_xB_z =
(\iu/\omega)(k_g^2 + \pi^2/a^2)\underline E = (\iu\omega/c^2)\underline E$:
$k_g^2 = \omega^2/c^2 - \pi^2/a^2$. (c) Brede wanden:
$\vect j_s = \pm(B_z\vect e_x - B_x\vect e_z)/\mu_0$ --- langs $z$ als
$\sin(\pi x/a)$, dwars als $\cos(\pi x/a)$ (nul in het midden); smalle
wanden: alleen langs $y$. (d) Een sleuf in het midden en in de
lengterichting snijdt alleen dwarse stromen door, die daar verdwijnen; een
dwarse sleuf snijdt de langsstromen door en straalt: de sleufantenne in een
golfgeleider.
\end{solution}

\begin{solution}{exo:b2:guided-waves:11}
(a)
$2a\sqrt{n_1^2 - n_2^2}/\lambda = 2 \times 50 \times 0.242/1.5 \approx 16$
per polarisatie. (b) $a < \lambda/2\mathrm{NA} = \qty{3.1}{\micro m}$. (c)
Een ronde kern met $\mathrm{NA} =
0.12$ is enkelmodig onder
$d = 2.405\lambda/\pi\mathrm{NA} = \qty{10}{\micro m}$: de standaard van
\qty{9}{\micro m}. (d) Een scherende straal wordt altijd volledig
teruggekaatst: de laagste mode bestaat bij elke frequentie, en haar
evanescente staarten spreiden zich alleen verder in de mantel uit.
\end{solution}

\begin{solution}{exo:b2:guided-waves:12}
(a) De straal buiten de as brengt haar tijd door in een lagere
brekingsindex, waar het licht sneller gaat: de langere weg wordt
gecompenseerd. (b) $n(r)\cos\theta = n_1\cos\theta_0$ met
$n \approx n_1(1 - \Delta r^2/a^2)$ en $\cos\theta \approx 1 - \theta^2/2$:
$(\dd r/\dd z)^2 =
\theta_0^2 - 2\Delta r^2/a^2$, een harmonische slingering
$r = r_0\sin(\sqrt{2\Delta}\,z/a)$. (c)
$2\pi a/\sqrt{2\Delta} = \qty{1.1}{mm}$. (d)
$n_1L\Delta^2/2c = \qty{0.25}{ns/km}$ tegen \qty{49}{ns/km}: tweehonderd
keer beter.
\end{solution}

\begin{solution}{pb:b2:guided-waves:1}
\textbf{1.} $c/2a = \qty{1.74}{GHz}$; TE$_{20}$ en TE$_{01}$:
\qty{3.49}{GHz}: bij \qty{2.45}{GHz} alleen TE$_{10}$.

\textbf{2.} $\sqrt{1 - (1.74/2.45)^2} = 0.70$: $\lambda_g = \qty{17.4}{cm}$,
$v_\varphi = 1.42c$, $v_g = 0.70c = \qty{2.1e8}{m/s}$.

\textbf{3.}
$E_0^2 = 4P/\varepsilon_0c\,ab\sqrt{\cdot} = 3200/6.9 \times 10^{-6}$:
$E_0 = \qty{21}{kV/m}$, een honderdste van de doorslag.

\textbf{4.} $j_s \approx E_0/\mu_0c = \qty{57}{A/m}$;
$\delta = \qty{10}{\micro m}$, $R_s = \qty{0.1}{\ohm}$;
$\tfrac12R_sj_s^2 \times 2a = \qty{28}{W/m}$: $3.5\%$ per meter --- staal is
verliesgevend, maar de geleider is kort.

\textbf{5.} $f = 1.5 \times 10^8\sqrt{11.1(m^2 + n^2) + 25p^2}$ binnen
$\pm3\%$ van \qty{2.45}{GHz}: $(1, 2, 3)$ en $(2, 1, 3)$ bij
\qty{2.51}{GHz}, $(4, 0, 2)$, $(0, 4, 2)$, $(4, 3, 0)$ en $(3, 4, 0)$ bij
\qty{2.50}{GHz} --- ongeveer zes.

\textbf{6.} Halve golflengten $a/m$, $b/n$, $d/p$: $7$ tot \qty{10}{cm} uit
elkaar. Het draaiplateau sleept het eten door maxima en minima; een menger
(een draaiende metalen waaier) schudt de modes door elkaar.

\textbf{7.} $f/Q = \qty{12}{MHz}$: de belaste resonanties overlappen tot een
continuüm; de magnetron vindt altijd een mode om te voeden.

\textbf{8.}
$W = QP/\omega = 3000 \times 800/1.54 \times 10^{10} = \qty{0.16}{J}$;
$u = \qty{8.7}{J/m^3}$;
$E_0 \approx \sqrt{4u/\varepsilon_0} = \qty{2}{MV/m}$ --- dicht bij de
doorslag: vlambogen; het teruggekaatste vermogen keert naar de magnetron
terug, die (een beetje) door een blindbelasting of een circulator wordt
beschermd --- vandaar ``laat hem nooit leeg draaien''.

\textbf{9.} $0.2/2.1 \times 10^8 = \qty{1}{ns}$.

\textbf{10.} \qty{1.6}{kW} piek; tussen de uitbarstingen galmt het veld uit
in $Q/\omega = \qty{13}{ns}$: de trilholte is het grootste deel van de tijd
leeg.

\textbf{11.} $\ell \approx a/\pi = \qty{0.48}{mm}$; $\eu^{-1/0.48}$:
\qty{-18}{dB} in amplitude, \qty{-36}{dB} in vermogen.

\textbf{12.} $\lambda/4 = \qty{3.1}{cm}$.

\textbf{13.} \qty{5}{mW/cm^2} $\times$ \qty{100}{cm^2} = \qty{0.5}{W}: $0.06\%$.

\textbf{14.} $f_1 \approx c/2d = \qty{3.75}{GHz} > \qty{2.45}{GHz}$: onder
de afsnijding; $\ell =
\qty{1.7}{cm}$, $\eu^{-30/1.7}$: \qty{-150}{dB}.

\textbf{15.} De wandstromen moeten ononderbroken in het gaas overlopen; een
onderbreking in de verbinding is een sleuf die hen doorsnijdt --- een
sleufantenne die naar buiten straalt.

\textbf{16.} $\mathrm{NA} = \sqrt{1.465^2 - 1.460^2} = 0.12$,
\qty{7}{\degree}: een smalle kegel, waarvoor een laser en een lens nodig
zijn.

\textbf{17.} $(n_1L/c)(n_1/n_2 - 1) = \qty{17}{ns/km}$; \qty{170}{ns} over
\qty{10}{km}: \qty{6}{Mbit/s}.

\textbf{18.} $\tfrac12(\pi d\,\mathrm{NA}/\lambda)^2 = 75$; voor
\qty{9}{\micro m}: $2.4$ --- het enkelmodige regime
($V = \pi d\,\mathrm{NA}/\lambda = 2.2 < 2.4$).

\textbf{19.} $\Delta x_0 = (c/n)\tau = \qty{1}{cm}$:
$t_{\text{sp}} = 10^{-4}/0.2 = \qty{0.5}{ms}$, \qty{100}{km}:
\qty{10}{Gbit/s} over die afstand, vóór compensatie.

\textbf{20.} \qty{20}{dB}, een factor $100$: \qty{10}{\micro W};
$6000/80 = 75$ versterkers.

\textbf{21.} $f = \qty{1.9e14}{Hz}$, $hf = \qty{0.8}{eV}$;
$8 \times 10^{15}$ fotonen per seconde; $8 \times 10^5$ per bit aan de
ingang, $8 \times 10^3$ na \qty{100}{km}.

\textbf{22.} De som van de twee verliezen is minimaal nabij
\qty{1.55}{\micro m} (\qty{0.2}{dB/km}).

\textbf{23.} $R = ((1.465 - 1)/2.465)^2 = 3.6\%$: \qty{0.16}{dB} per vlak,
twee keer bij een connector plus de interferentie van de spleet; gel of
fysiek contact verwijdert de lucht.

\textbf{24.} Bij een zachte bocht ontmoet de straal de rand nog altijd
voorbij de grenshoek; bij een scherpe zakt de inval op de buitenwand eronder
en lekt er licht weg.

\textbf{25.} Oven: metalen wanden, reflectie met $r = -1$; de bandbreedte
wordt door de modes van de trilholte bepaald; de verliezen zitten in de huid
van het staal. Vezel: \omterm{thm:b2:wave-interfaces:snell}{totale interne weerkaatsing} aan een grensvlak
glas--glas; de bandbreedte wordt door modale en chromatische dispersie
bepaald; de verliezen komen van verstrooiing en absorptie, \qty{0.2}{dB/km}.
\end{solution}
